\documentclass[border=10pt]{standalone}
\usepackage{fontspec}
\setmainfont{Noto Serif}
\newfontfamily\bengalifont{Noto Serif Bengali}[Script=Bengali]
\usepackage[Latin,Bengali]{ucharclasses}
\setTransitionsFor{Bengali}{\begingroup\bengalifont}{\endgroup}
\usepackage{newunicodechar}
\newunicodechar{।}{{\bengalifont\char"0964}}
\newunicodechar{॥}{{\bengalifont\char"0965}}

\usepackage[table]{xcolor}
\usepackage{booktabs,tabularx,array}
\definecolor{bapilrule}{HTML}{1F2937}
\definecolor{bapilstripe}{HTML}{F3F4F6}
\renewcommand{\arraystretch}{1.28}
\setlength{\tabcolsep}{5pt}

\begin{document}
\begin{minipage}{17cm}
{\large\bfseries সারণি ৫১।}\ {\large দিগন্তভিত্তিক নির্দেশক অর্থায়ন-মিশ্রণ}\par\vspace{7pt}
\small
\rowcolors{2}{bapilstripe}{white}
\begin{tabularx}{\linewidth}{@{}>{\hsize=1.512\hsize\raggedright\arraybackslash}X>{\hsize=0.872\hsize\raggedright\arraybackslash}X>{\hsize=0.872\hsize\raggedright\arraybackslash}X>{\hsize=0.872\hsize\raggedright\arraybackslash}X>{\hsize=0.872\hsize\raggedright\arraybackslash}X@{}}
\arrayrulecolor{bapilrule}\toprule
\rowcolor{white}\textbf{উৎস} & \textbf{H1 ২০২৭ থেকে ২০২৮} & \textbf{H2 ২০২৯ থেকে ২০৩০} & \textbf{H3 ২০৩১ থেকে ২০৩৫} & \textbf{H4 ২০৩৬ থেকে ২০৪১} \\
\midrule
সরকারের মূল অবদান (সমঝোতা স্মারক, উন্নয়ন বাজেট) & ৩০ থেকে ৪০\% & ৪০ থেকে ৫০\% & ৪৫ থেকে ৫৫\% & ৫০ থেকে ৬০\% \\
প্রতিযোগিতামূলক গবেষণা অনুদান & ২০ থেকে ৩০\% & ১৫ থেকে ২৫\% & ১০ থেকে ২০\% & ১০ থেকে ১৫\% \\
উন্নয়ন সহযোগী (পাইলটভিত্তিক) & ২০ থেকে ৩০\% & ১৫ থেকে ২০\% & ৫ থেকে ১০\% & ০ থেকে ৫\% \\
সংস্থার সঙ্গে সেবা-চুক্তি (ব্যয় পুনরুদ্ধার) & ০\% & ৫ থেকে ১০\% & ১৫ থেকে ২৫\% & ২০ থেকে ৩০\% \\
শিল্প ও জনহিতৈষী সহায়তা (উন্মুক্ত বেঞ্চমার্ক, ফেলোশিপ) & ৫ থেকে ১০\% & ৫ থেকে ১০\% & ৫ থেকে ১০\% & ৫ থেকে ১০\% \\
\bottomrule
\end{tabularx}
\par\vspace{5pt}{\footnotesize\itshape কোনো বছরে সরকার-বহির্ভূত কোনো একক অর্থদাতা বাজেটের ৩০\%-এর বেশি দেবেন না, এবং শিল্পখাতের সহায়তা কখনো গবেষণার ফলাফল বা শাসনগত সিদ্ধান্তে প্রভাব কিনতে পারবে না।}
\par\vspace{3pt}{\scriptsize\color{gray} উৎস: BAPIL কারিগরি ব্লুপ্রিন্ট, সংস্করণ ১.০ (২০২৬)। CC BY 4.0।}
\end{minipage}
\end{document}
